% Provenance (never printed): see sources/outline_sources.md (Year 2: radiative forcing
% quantified; feedbacks and climate sensitivity; detection and attribution; scenarios; the
% remaining carbon budget; components of sea-level rise; physical science only). Every statement
% from IPCC AR6 WG1 (SPM; chapters 3, 5, 7 and 9; Annex III data), ledger rows in
% sources/ledger/book2.md (ipcc:*) and sources/ledger/book4.md (ar6*:*). Curves and numbers from
% figdata/bachelor-2/anthropogenic.py and the Book 4 second-writer scratch calculator calc/ch27.py.
\chapter{A mudança climática antrópica}\label{ch:b2:anthropogenic-climate-change}

Desde 1970 a superfície do globo aqueceu mais depressa que em qualquer outro período de cinquenta anos, ao menos nos
últimos dois mil anos (\emph{alta confiança}). Os volumes escolares contaram o que mudou e
por quê. Este capítulo transforma esse relato num balanço: quantos watts por metro quadrado os fatores
humanos acrescentam, quantos o planeta mais quente devolve ao espaço, quanto aquecimento cada tonelada de dióxido
de carbono produz, e como a subida do mar se reparte entre a água que aquece e o gelo que derrete. Cada número
vem com sua incerteza, na linguagem calibrada das avaliações.
% ledger: ipcc:fastest-2000yr, ipcc:likelihood

\begin{recall}
\textbf{Dos volumes escolares}: o forçamento radiativo, $\qty{3.93}{W.m^{-2}}$ por duplicação do
dióxido de carbono, e sua lei logarítmica; a retroalimentação climática; a \omterm{def:b2:anthropogenic-climate-change:ecs}{sensibilidade climática de equilíbrio}, melhor
estimativa \qty{3}{\celsius}, \emph{provável} 2.5 a \qty{4}{\celsius}; a \omterm{def:b2:anthropogenic-climate-change:da}{atribuição}, o cenário de emissões
e a projeção; a linguagem calibrada: \emph{provável} 66 a \qty{100}{\percent}, \emph{muito
provável} 90 a \qty{100}{\percent}, um intervalo $[x \text{ a } y]$ sendo um intervalo muito provável. A influência
humana aqueceu o clima em \qty{1.07}{\celsius} entre 1850--1900 e 2010--2019.
\textbf{Dos livros de física}: o corpo negro e a lei de Stefan--Boltzmann; a dilatação térmica.
\textbf{Dos livros de biologia}: o ciclo do carbono, a fração que permanece no ar, os sumidouros de carbono e o
balanço de carbono da biosfera. \textbf{Do \cref{ch:b2:deep-time-climate}}: o ciclo lento do carbono,
que só retira um excesso de dióxido de carbono ao longo de centenas de milhares de anos.
\end{recall}
% ledger: ipcc:erf-2xco2, ipcc:log-co2, ipcc:ecs, ipcc:likelihood, ipcc:attribution

\section{O forçamento quantificado}\label{sec:b2:anthropogenic-climate-change:forcing}

\begin{definition}[Forçamento radiativo efetivo]\label{def:b2:anthropogenic-climate-change:erf}
O \emph{forçamento radiativo efetivo}\index{forçamento radiativo efetivo} (ERF) de uma perturbação
--- uma mudança num gás de efeito estufa, nos aerossóis, no Sol --- é a variação que ela causa no fluxo
líquido descendente de radiação no topo da atmosfera, em \unit{W.m^{-2}}, depois que a atmosfera,
as nuvens e a superfície terrestre se ajustaram a ela, mas antes que a temperatura média da superfície do globo
tenha respondido. Os ajustes contados são os causados pelo próprio agente, independentemente da
temperatura da superfície.
\end{definition}
% ledger: ar6ch7:erf-def

A definição traça a linha entre forçamento e resposta pelo que depende da temperatura da superfície,
não pela rapidez. O forçamento instantâneo do dióxido de carbono é a variação de fluxo logo
depois que o gás é acrescentado; a estratosfera, que então esfria em alguns meses, altera-o uma primeira vez;
a troposfera, seu vapor d'água e suas nuvens alteram-no de novo. Para uma duplicação do dióxido de
carbono a avaliação encontra \qty{3.75}{W.m^{-2}} depois do ajuste da estratosfera, e
ajustes na troposfera que se compensam em parte: \qty{-0.60}{W.m^{-2}} por seu aquecimento, +0.22
pelo vapor d'água, +0.45 pelas nuvens e +0.11 pela superfície terrestre, \qty{+0.18}{W.m^{-2}} no total.
O ERF vale $3.75 + 0.18 = \qty{3.93}{W.m^{-2}}$, \qty{6}{\percent} a mais que os
\qty{3.71}{W.m^{-2}} da avaliação anterior (\cref{exo:b2:anthropogenic-climate-change:1}).
% ledger: ar6ch7:2xco2, ar6ch7:erf-def

O forçamento de cada agente desde 1750 foi reconstruído ano a ano. A figura abaixo o mostra:
o dióxido de carbono a \qty{2.16}{W.m^{-2}} em 2019, os outros gases de efeito estufa --- metano,
óxido nitroso, gases halogenados e ozônio --- a \qty{1.69}{W.m^{-2}} juntos, e os
aerossóis a \qty{-1.06}{W.m^{-2}}. O forçamento humano total é de \qty{2.72}{W.m^{-2}}, e seu intervalo muito
provável, 1.96 a \qty{3.48}{W.m^{-2}}, é largo sobretudo por causa dos aerossóis. Seu ERF em
1750--2014 é avaliado em $-1.3$ [$-2.0$ a $-0.6$] \unit{W.m^{-2}} (\emph{média confiança}),
em sua maior parte pelo efeito sobre as nuvens, que ficam mais brilhantes quando os aerossóis fornecem mais gotículas.
As erupções vulcânicas acrescentam picos negativos breves (\cref{ch:b2:eruptive-styles}): o mais profundo do
registro, \qty{-4.8}{W.m^{-2}} em 1816, segue a erupção do Tambora em 1815.
% ledger: data:b4-ar6-erf, ipcc:erf-total, ipcc:erf-agents, ar6ch7:aerosol, usgs:msh-comparisons

\begin{omfigure}
\centering
\begin{tikzpicture}
\begin{axis}[omrecord, width=0.92\linewidth, height=6.2cm, xmin=1750, xmax=2020, ymin=-5.2,
  ymax=3.8, xlabel={ano}, ylabel={ERF desde 1750 (\unit{W.m^{-2}})},
  xtick={1750,1800,...,2000}, /pgf/number format/1000 sep={}, xtick pos=bottom, ytick pos=left,
  tick label style={font=\scriptsize}, label style={font=\small},
  legend style={at={(0.5,-0.24)}, anchor=north, font=\scriptsize, draw=black!30,
  legend columns=3, /tikz/every even column/.append style={column sep=6pt}},
  legend cell align=left]
% computed in figdata/bachelor-2/anthropogenic.py part erf; ledger: data:b4-ar6-erf
\addplot[name path=hi, draw=none, forget plot] table[x=year, y=hi] {figdata/out/bachelor-2-anthropogenic-erf.dat};
\addplot[name path=lo, draw=none, forget plot] table[x=year, y=lo] {figdata/out/bachelor-2-anthropogenic-erf.dat};
\addplot[black!15, forget plot] fill between[of=hi and lo];
\addplot[black!45, line width=0.4pt] table[x=year, y=natural] {figdata/out/bachelor-2-anthropogenic-erf.dat};
\addlegendentry{natural (vulcões, Sol)}
\addplot[omThm, line width=1pt] table[x=year, y=co2] {figdata/out/bachelor-2-anthropogenic-erf.dat};
\addlegendentry{dióxido de carbono}
\addplot[omThm!50!orange, line width=1pt, densely dashed] table[x=year, y=other_ghg] {figdata/out/bachelor-2-anthropogenic-erf.dat};
\addlegendentry{outros gases de efeito estufa e ozônio}
\addplot[omDef, line width=1pt] table[x=year, y=aerosol] {figdata/out/bachelor-2-anthropogenic-erf.dat};
\addlegendentry{aerossóis}
\addplot[black, line width=1.4pt] table[x=year, y=anthro] {figdata/out/bachelor-2-anthropogenic-erf.dat};
\addlegendentry{total humano (cinza: 5--95\,\%)}
\draw[black!40] (axis cs:1750,0) -- (axis cs:2020,0);
\node[anchor=west, font=\scriptsize, text=black!60, inner sep=1pt] at (axis cs:1817.5,-4.75) {$\leftarrow$ Tambora};
\end{axis}
\end{tikzpicture}
\par\smallskip{\small \omterm{def:b2:anthropogenic-climate-change:erf}{Forçamento radiativo efetivo} de 1750 a 2019, em relação a 1750, por fator.
O forçamento humano total atinge \qty{2.72}{W.m^{-2}} em 2019, com o intervalo muito provável sombreado.
Dados: Sexto Relatório de Avaliação do IPCC, Grupo de Trabalho~I, dados estendidos do Anexo~III (CC BY 4.0).}
\end{omfigure}
% values: anthro 2.7204 (1.9569-3.4819), co2 2.1563, other_ghg 1.6850, aerosol -1.0583, volcanic -4.8162 in 1816 (anthropogenic.py numbers)

\section{O balanço de energia}\label{sec:b2:anthropogenic-climate-change:energy}

Um forçamento positivo faz o planeta ganhar energia. À medida que aquece, ele irradia mais para o espaço, e o
ganho diminui; o clima só se estabiliza quando os dois se equilibram.

\begin{definition}[Desequilíbrio energético da Terra; parâmetro de retroalimentação climática]\label{def:b2:anthropogenic-climate-change:eei}
O \emph{desequilíbrio energético da Terra}\index{desequilíbrio energético da Terra} $N$ é o fluxo líquido de energia que entra
no sistema climático no topo da atmosfera, em média sobre o globo e sobre um ano ou mais,
em \unit{W.m^{-2}}: a luz solar absorvida menos o infravermelho emitido. É a taxa à qual o oceano,
os continentes, o gelo e o ar ganham calor, por metro quadrado da superfície da Terra.
\end{definition}

\begin{definition}[Parâmetro de retroalimentação climática]\label{def:b2:anthropogenic-climate-change:lambda}
O \emph{parâmetro de retroalimentação climática}\index{parâmetro de retroalimentação climática} $\lambda$ é a variação do
fluxo líquido descendente no topo da atmosfera por grau de aquecimento da superfície do globo, em
\unit{W.m^{-2}.\celsius^{-1}}, causada pela resposta do sistema climático a esse aquecimento. Um
clima estável tem $\lambda < 0$: o aquecimento aumenta a perda de energia para o espaço.
\end{definition}

O volume escolar escreveu o balanço com o módulo de $\lambda$, um número positivo; daqui em diante o
sinal é mantido, como nas avaliações, que escrevem o parâmetro $\alpha$.
% ledger: ar6ch7:energy-budget

\begin{proposition}[Balanço global de energia]\label{prop:b2:anthropogenic-climate-change:balance}
Seja $\Delta F$ o \omterm{def:b2:anthropogenic-climate-change:erf}{forçamento radiativo efetivo} e $\Delta T$ o aquecimento da superfície do globo
em relação a um estado de equilíbrio. Se a resposta do fluxo emitido é linear em $\Delta T$,
\[
  \Delta N = \Delta F + \lambda\,\Delta T .
\]
Num novo equilíbrio $\Delta N = 0$, logo $\Delta T_{\text{eq}} = -\Delta F/\lambda$.
\end{proposition}
\begin{proof}
Por definição o forçamento altera o fluxo líquido descendente em $\Delta F$ com temperatura da superfície
fixa, e a resposta ao aquecimento altera-o em $\lambda\,\Delta T$ em primeira ordem;
o desequilíbrio é a soma dos dois. Igualá-lo a zero dá o equilíbrio.
\end{proof}

Medir $N$ significa medir o calor que o sistema ganha, e o oceano retém quase todo ele.
De 1971 a 2018 o sistema climático ganhou \qty{435}{ZJ} (zettajoules, \qty{e21}{J}): o oceano
\qty{91}{\percent}, \qty{242}{ZJ} em seus \qty{700}{m} superiores e \qty{154}{ZJ} abaixo; os continentes
\qty{5}{\percent}, o gelo que derrete \qty{3}{\percent}, a atmosfera \qty{1}{\percent}. O desequilíbrio
médio foi de \qty{0.57}{W.m^{-2}} em 1971--2018, e aumentou: \qty{0.79}{W.m^{-2}}
[0.52 a 1.06] em 2006--2018 (\cref{exo:b2:anthropogenic-climate-change:2}).
% ledger: ar6ch7:inventory, ipcc:heating

\begin{omfigure}
\centering
\begin{tikzpicture}[font=\small, >=Stealth]
\draw[fill=omDef!8, draw=omDef!60] (-3.2,0) rectangle (3.2,1.6);
\node[align=center] at (0,1.05) {sistema climático};
\node[align=center, font=\scriptsize] at (0,0.5) {oceano \qty{91}{\percent}, continentes \qty{5}{\percent}, gelo \qty{3}{\percent}, ar \qty{1}{\percent}};
\draw[->, line width=2.2pt, omThm] (-2.0,3.3) -- node[pos=0.12, left=3pt, align=right, font=\scriptsize] {forçamento $\Delta F$} (-2.0,1.7);
\draw[->, line width=2.2pt, omDef] (2.0,1.7) -- node[pos=0.88, right=3pt, align=left, font=\scriptsize] {resposta $\lambda\,\Delta T$} (2.0,3.3);
\draw[->, line width=1.4pt, black!70] (0,-0.1) -- node[right=3pt, align=left, font=\scriptsize] {armazenado: $\Delta N = \Delta F + \lambda\,\Delta T$} (0,-1.0);
\draw[black!40, dashed] (-3.4,2.4) -- (3.4,2.4) node[right, font=\scriptsize, text=black!60] {topo da atmosfera};
\end{tikzpicture}
\par\smallskip{\small O balanço de energia de um clima forçado (esquemático). O forçamento acrescenta energia no
topo da atmosfera, a resposta do planeta mais quente retira uma parte ($\lambda\,\Delta T$ é
negativo), e a diferença é armazenada, sobretudo no oceano. Partes do calor ganho de 1971
a 2018.}
\end{omfigure}
% ledger: ar6ch7:inventory

\begin{method}[Uma sensibilidade climática a partir do balanço de energia histórico]\label{met:b2:anthropogenic-climate-change:energy-budget}
\begin{enumerate}
  \item Escolha dois períodos, uma referência perto do equilíbrio e um recente; tome a variação do
        forçamento $\Delta F$, do desequilíbrio $\Delta N$ e da temperatura global $\Delta T$.
  \item Resolva o balanço para o parâmetro de retroalimentação: $\lambda = (\Delta N - \Delta F)/\Delta T$.
  \item A sensibilidade a uma duplicação é então $-F_{2\times}/\lambda$; o aquecimento no momento do
        forçamento, escalonado para $F_{2\times}$, é $F_{2\times}\Delta T/\Delta F$.
\end{enumerate}
\end{method}

De 1850--1900 a 2006--2019 a avaliação toma $\Delta N = \qty{0.59}{W.m^{-2}}$, $\Delta T =
\qty{1.03}{\celsius}$ e $\Delta F = \qty{2.20}{W.m^{-2}}$. Então $\lambda = (0.59 - 2.20)/1.03 =
\qty{-1.56}{W.m^{-2}.\celsius^{-1}}$, e a sensibilidade é $3.93/1.56 = \qty{2.5}{\celsius}$,
o valor que a avaliação encontra, com um intervalo muito provável de 1.6 a \qty{4.8}{\celsius}. Mas esta é
uma sensibilidade \emph{efetiva}, para as retroalimentações do aquecimento até agora. O aquecimento foi
mais forte onde as retroalimentações são mais negativas, no Pacífico tropical ocidental, e mais fraco no
Pacífico oriental e no oceano Austral; à medida que esse padrão evolui, espera-se que o parâmetro de retroalimentação
se torne menos negativo (\emph{alta confiança}), de modo que a sensibilidade de longo prazo é maior que a
estimativa histórica (\cref{exo:b2:anthropogenic-climate-change:7}).
% ledger: ar6ch7:energy-budget, ar6ch7:pattern, ipcc:erf-2xco2
% values: lambda_energy -1.563, ecs_energy 2.514 (anthropogenic.py numbers)

\section{Retroalimentações e sensibilidade}\label{sec:b2:anthropogenic-climate-change:feedbacks}

\begin{definition}[Resposta de Planck; retroalimentações do vapor d'água, do gradiente vertical e das nuvens]\label{def:b2:anthropogenic-climate-change:feedbacks}
O parâmetro de retroalimentação é dividido nas contribuições das grandezas que mudam com a
temperatura da superfície. A \emph{resposta de Planck}\index{resposta de Planck} é o infravermelho adicional
emitido para o espaço quando a superfície e toda a troposfera aquecem da mesma quantidade. A
\emph{retroalimentação do vapor d'água}\index{retroalimentação do vapor d'água} é a variação de fluxo devida ao aumento
do vapor d'água atmosférico com a temperatura. A \emph{retroalimentação do gradiente
vertical}\index{retroalimentação do gradiente vertical} é a variação de fluxo devida a um aquecimento que não é uniforme na vertical. A \emph{retroalimentação
das nuvens}\index{retroalimentação das nuvens} é a variação de fluxo devida à mudança da quantidade, da altitude e do
brilho das nuvens.
\end{definition}

\begin{proposition}[A resposta de Planck]\label{prop:b2:anthropogenic-climate-change:planck}
Se a Terra emite para o espaço uma fração fixa $1 - g$ do infravermelho $\sigma T^4$ de sua superfície à
temperatura $T$, sua \omterm{def:b2:anthropogenic-climate-change:feedbacks}{resposta de Planck} é
\[
  \lambda_{\text P} = -(1 - g)\,4\sigma T^3 .
\]
Com $g = 0.4$ e $T = \qty{288}{K}$, $\lambda_{\text P} = \qty{-3.25}{W.m^{-2}.K^{-1}}$.
\end{proposition}
\begin{proof}
O fluxo emitido é $(1 - g)\sigma T^4$; sua derivada em relação a $T$ é $(1 - g)\,4\sigma
T^3$, e uma perda adicional é uma contribuição negativa ao fluxo líquido descendente. Numericamente $0.6
\times 4 \times \num{5.670e-8} \times 288^3 = 3.25$.
\end{proof}
% ledger: ar6ch7:planck, const:sigma, ipcc:budget-surface
% values: planck -3.2509 (anthropogenic.py numbers)

A fração $g$ é o efeito estufa normalizado: a atmosfera retém \qty{159}{W.m^{-2}} dos
\qty{398}{W.m^{-2}} que a superfície emite, $g \approx 0.4$. A avaliação, com códigos de radiação
detalhados, encontra $-3.22$ [$-3.4$ a $-3.0$] (\emph{alta confiança}). O vapor d'água segue
a lei de Clausius--Clapeyron em escala global, e mais vapor absorve mais infravermelho: a
\omterm{def:b2:anthropogenic-climate-change:feedbacks}{retroalimentação do vapor d'água} é fortemente positiva, \qty{1.85}{W.m^{-2}.\celsius^{-1}} segundo os satélites.
Nos trópicos a alta troposfera aquece mais que a superfície, e então irradia mais: a
\omterm{def:b2:anthropogenic-climate-change:feedbacks}{retroalimentação do gradiente vertical} é negativa ali, positiva perto dos polos, e de cerca de $-0.50$ globalmente. As duas
são avaliadas juntas, \num{1.30} [1.1 a 1.5], porque o que um modelo põe numa tende a
tirar da outra. A neve e o gelo que derretem escurecem a superfície, +0.35; as nuvens acrescentam \num{0.42}
[$-0.10$ a $0.94$], a maior incerteza. Uma retroalimentação líquida negativa das nuvens é \emph{muito improvável}.
% ledger: ar6ch7:planck, ar6ch7:wv-lr, ar6ch7:table710, ipcc:cloud-negative, ipcc:budget-surface

\begin{omfigure}
\centering
\begin{tikzpicture}
\begin{axis}[xbar, width=0.78\linewidth, height=5.4cm, xmin=-3.7, xmax=1.8, ymin=0.4, ymax=6.6,
  ytick={1,...,6}, yticklabels={líquida, biogeofísicas e químicas, nuvens, albedo da superfície,
  vapor d'água + gradiente vertical, Planck}, y tick label style={font=\scriptsize},
  tick label style={font=\scriptsize}, xlabel={parâmetro de retroalimentação (\unit{W.m^{-2}.\celsius^{-1}})},
  label style={font=\small}, bar width=7pt, axis x line=bottom, axis y line=left,
  axis line style={-}, xmajorgrids, grid style={gray!22, line width=0.3pt}, extra x ticks={0},
  extra x tick style={grid style={black!60, line width=0.5pt}}]
% ledger: ar6ch7:table710 (central values and very likely ranges)
\addplot[fill=omDef!55, draw=omDef!80!black, error bars/.cd, x dir=both, x explicit,
  error bar style={black!70}] coordinates {
  (-3.22,6) += (0.22,0) -= (0.18,0)
  (-0.01,2) += (0.26,0) -= (0.26,0)
  (-1.16,1) += (0.65,0) -= (0.65,0)};
\addplot[fill=omThm!55, draw=omThm!80!black, error bars/.cd, x dir=both, x explicit,
  error bar style={black!70}] coordinates {
  (1.30,5) += (0.20,0) -= (0.20,0)
  (0.35,4) += (0.25,0) -= (0.25,0)
  (0.42,3) += (0.52,0) -= (0.52,0)};
\end{axis}
\end{tikzpicture}
\par\smallskip{\small As retroalimentações climáticas avaliadas pelo IPCC (Sexto Relatório de Avaliação, Grupo
de Trabalho~I, Tabela 7.10), valores centrais e intervalos muito prováveis. As retroalimentações positivas amplificam um aquecimento
(vermelho), as negativas o amortecem (azul); a líquida é a soma.}
\end{omfigure}

\begin{definition}[Sensibilidade climática de equilíbrio]\label{def:b2:anthropogenic-climate-change:ecs}
A \emph{sensibilidade climática de equilíbrio}\index{sensibilidade climática de equilíbrio} (ECS) é o
aquecimento da superfície do globo, no novo equilíbrio, causado por uma duplicação do dióxido de carbono do
ar a partir de seu nível pré-industrial, mantendo-se fixos os mantos de gelo.
\end{definition}

\begin{proposition}[Sensibilidade e retroalimentação]\label{prop:b2:anthropogenic-climate-change:ecs}
Se o parâmetro de retroalimentação é constante, $\text{ECS} = -F_{2\times}/\lambda$. Com $F_{2\times} =
\qty{3.93}{W.m^{-2}}$ e o valor líquido $\lambda = \qty{-1.16}{W.m^{-2}.\celsius^{-1}}$ das retroalimentações
acima, $\text{ECS} = \qty{3.4}{\celsius}$; a \omterm{def:b2:anthropogenic-climate-change:feedbacks}{resposta de Planck} sozinha daria $3.93/3.22 =
\qty{1.2}{\celsius}$.
\end{proposition}
\begin{proof}
Aplique-se a \cref{prop:b2:anthropogenic-climate-change:balance} no equilíbrio com $\Delta F =
F_{2\times}$.
\end{proof}
% ledger: ar6ch7:table710, ipcc:erf-2xco2
% values: ecs_process 3.388 (anthropogenic.py numbers)

As retroalimentações são uma linha de evidência; o aquecimento desde 1850, o balanço de energia e os climas
do passado são outras. Combinadas, dão a ECS avaliada de \qty{3}{\celsius}, \emph{provável}
2.5 a \qty{4}{\celsius} (\emph{alta confiança}), \emph{muito provável} 2 a \qty{5}{\celsius}.
Como a ECS varia como $1/\lambda$, uma incerteza simétrica sobre $\lambda$ dá um intervalo
de ECS esticado em direção aos valores altos (\cref{exo:b2:anthropogenic-climate-change:12}).
% ledger: ipcc:ecs

\begin{definition}[Resposta climática transitória]\label{def:b2:anthropogenic-climate-change:tcr}
A \emph{resposta climática transitória}\index{resposta climática transitória} (TCR) é o aquecimento da superfície
do globo no momento em que o dióxido de carbono do ar duplicou, quando ele sobe \qty{1}{\percent}
por ano a partir de seu nível pré-industrial, cerca de setenta anos depois do início.
\end{definition}

A TCR é menor que a ECS porque o oceano ainda está absorvendo calor quando a duplicação é
atingida. Sua melhor estimativa é \qty{1.8}{\celsius}, \emph{provável} 1.4 a \qty{2.2}{\celsius}
(\emph{alta confiança}); o balanço de energia histórico sozinho dá $3.93 \times 1.03/2.20 =
\qty{1.8}{\celsius}$ também.
% ledger: ar6ch7:tcr, ar6ch7:energy-budget
% values: tcr_energy 1.840 (anthropogenic.py numbers)

\section{Detecção e atribuição}\label{sec:b2:anthropogenic-climate-change:attribution}

\begin{definition}[Detecção; atribuição; impressão digital climática]\label{def:b2:anthropogenic-climate-change:da}
A \emph{detecção}\index{detecção (da mudança climática)} de uma mudança climática consiste em mostrar que o
clima mudou num sentido estatístico definido, além do que sua variabilidade interna poderia
produzir, sem indicar a causa. A \emph{atribuição}\index{atribuição} consiste em avaliar as
contribuições de várias causas para uma mudança detectada, com uma probabilidade atribuída. Uma \emph{impressão
digital climática}\index{impressão digital climática} é o padrão no espaço e no tempo da resposta a um
fator, simulado por modelos climáticos, pelo qual esse fator é reconhecido nas observações.
\end{definition}

Fatores diferentes deixam impressões digitais diferentes. Mais gás de efeito estufa aquece a troposfera e esfria
a estratosfera, porque um cobertor infravermelho mais espesso faz as camadas superiores irradiarem mais para o espaço;
um Sol mais brilhante aqueceria as duas. As observações mostram o primeiro padrão: a troposfera aqueceu
desde pelo menos a década de 1950, e é \emph{virtualmente certo} que a estratosfera esfriou. É
\emph{muito provável} que os gases de efeito estufa humanos tenham sido o principal fator do aquecimento troposférico
desde 1979, e \emph{extremamente provável} que a destruição da camada de ozônio tenha sido o principal fator do
resfriamento da baixa estratosfera de 1979 a meados da década de 1990.
% ledger: ar6ch3:upper-air

\begin{method}[A impressão digital ótima]\label{met:b2:anthropogenic-climate-change:fingerprinting}
\begin{enumerate}
  \item Simule a resposta a cada fator separadamente: as impressões digitais $x_1, \dots, x_k$, como
        vetores de valores no espaço e no tempo.
  \item Escreva as observações como $y = \sum_i \beta_i x_i + \varepsilon$, com $\varepsilon$ a
        variabilidade interna, cujas estatísticas são tiradas de longas simulações sem forçamento.
  \item Estime os fatores de escala $\beta_i$ por mínimos quadrados, ponderados pelo inverso da
        variabilidade para que as componentes mais ruidosas contem menos, com seus intervalos de confiança.
  \item Um $\beta_i$ significativamente maior que zero significa que a resposta ao fator $i$ é
        detectada; um intervalo que contém 1 e exclui 0 significa que a resposta simulada é
        coerente com as observações, e a mudança pode ser atribuída em parte a esse fator.
\end{enumerate}
\end{method}
% ledger: ar6ch3:fingerprint

Aplicado à temperatura da superfície, o método dá a \omterm{def:b2:anthropogenic-climate-change:da}{atribuição} lembrada no início: dos
\qty{1.06}{\celsius} observados entre 1850--1900 e 2010--2019, a influência humana \emph{provavelmente}
causou 0.8 a \qty{1.3}{\celsius}, melhor estimativa \qty{1.07}{\celsius}; os fatores naturais e a
variabilidade interna contribuem com não mais que alguns décimos de grau num sentido ou no outro
(\cref{exo:b2:anthropogenic-climate-change:9,exo:b2:anthropogenic-climate-change:10}).
% ledger: ipcc:attribution, ipcc:observed-2010-2019

\section{Cenários e emissões cumulativas}\label{sec:b2:anthropogenic-climate-change:scenarios}

\begin{definition}[Trajetória Socioeconômica Compartilhada]\label{def:b2:anthropogenic-climate-change:ssp}
Uma \emph{Trajetória Socioeconômica Compartilhada}\index{Trajetória Socioeconômica Compartilhada} (SSP) é uma de cinco
narrativas de como a população, as economias e a tecnologia poderiam evoluir ao longo do século. Os
cenários de emissões da avaliação são rotulados SSP$x$-$y$: $x$ designa a trajetória, e $y$ o
forçamento radiativo aproximado, em \unit{W.m^{-2}}, que o cenário atinge em 2100.
\end{definition}
% ledger: ar6spm:ssp

Os cinco cenários principais vão de SSP1-1.9 a SSP5-8.5. A avaliação da ciência física toma-os como
dados: não pondera sua probabilidade nem avalia as narrativas. Para cada um, os modelos climáticos
dão um aquecimento projetado para 2081--2100 acima de 1850--1900, de \qty{1.4}{\celsius} (muito provável 1.0
a 1.8) no SSP1-1.9 a \qty{4.4}{\celsius} (3.3 a 5.7) no SSP5-8.5. As projeções compartilham uma
lei simples.
% ledger: ar6spm:ssp, ipcc:proj-temp

\begin{definition}[Resposta climática transitória às emissões cumulativas de CO$_2$]\label{def:b2:anthropogenic-climate-change:tcre}
A \emph{resposta climática transitória às emissões cumulativas de \ce{CO2}}\index{resposta climática transitória às emissões cumulativas de CO$_2$}
(TCRE) é o aquecimento da superfície do globo por unidade de dióxido de carbono
emitido no total desde o período pré-industrial, em \unit{\celsius} por
\qty{1000}{Gt} de \ce{CO2}.
\end{definition}

\begin{proposition}[O aquecimento segue as emissões cumulativas]\label{prop:b2:anthropogenic-climate-change:tcre}
Com boa aproximação o aquecimento causado pelo dióxido de carbono é proporcional ao total emitido,
qualquer que seja a trajetória das emissões: $\Delta T \approx \text{TCRE} \times E$, com $E$ as emissões
cumulativas.
\end{proposition}
\admitted
% ledger: ar6spm:tcre, ar6ch5:tcre-linear

A lei não é óbvia. O forçamento cresce apenas como o logaritmo da concentração, de modo que cada
tonelada acrescentada aquece menos; mas o oceano absorve calor e carbono com cada vez menos eficiência à medida que
as emissões se acumulam, de modo que mais de cada tonelada fica no ar e seu aquecimento é menos retardado. Os
dois efeitos quase se compensam, o que todos os estudos concordam ser a origem da lei. A relação é
quase linear (\emph{alta confiança}), e cada \qty{1000}{Gt} de dióxido de carbono \emph{provavelmente} causa
0.27 a \qty{0.63}{\celsius} de aquecimento, melhor estimativa \qty{0.45}{\celsius}, ou \qty{1.65}{\celsius}
por \qty{1000}{PgC} de carbono (\cref{exo:b2:anthropogenic-climate-change:4}). De 1850 a 2019,
foram emitidas $\num{2390} \pm \qty{240}{Gt}$ de dióxido de carbono.
% ledger: ar6ch5:tcre-linear, ar6spm:tcre, ar6spm:budget

\begin{definition}[Orçamento de carbono restante]\label{def:b2:anthropogenic-climate-change:budget}
O \emph{orçamento de carbono restante}\index{orçamento de carbono restante} para um limite de temperatura é o
dióxido de carbono cumulativo que ainda pode ser emitido, a partir de uma data indicada, com o aquecimento global ficando
abaixo desse limite com uma probabilidade indicada, levando-se em conta o aquecimento causado pelos outros fatores
humanos.
\end{definition}

\begin{method}[O orçamento de carbono restante a partir da TCRE]\label{met:b2:anthropogenic-climate-change:budget}
\begin{enumerate}
  \item Subtraia do limite o aquecimento humano já atingido,
        \qty{1.07}{\celsius} para 2010--2019.
  \item Subtraia o aquecimento que os outros fatores terão acrescentado quando as emissões de dióxido de carbono
        chegarem a zero líquido.
  \item Divida pela TCRE: este é o orçamento a partir do meio de 2010--2019.
  \item Subtraia as emissões desde então, e uma margem para as retroalimentações lentas do sistema terrestre que os
        modelos deixam de fora (o degelo do permafrost é a maior).
\end{enumerate}
\end{method}

A avaliação combina cinco componentes desta maneira: a TCRE; o aquecimento histórico; um
aquecimento adicional pelos outros fatores de cerca de 0.1 a \qty{0.2}{\celsius}; o aquecimento que pode
continuar depois que as emissões param, avaliado em zero em média (\emph{provavelmente} dentro de $\pm
\qty{0.19}{\celsius}$); e as retroalimentações do sistema terrestre, cerca de \qty{26}{Gt} de dióxido de carbono por
grau. Para \qty{1.5}{\celsius}, \qty{0.43}{\celsius} acima de 2010--2019, o orçamento restante a partir do
início de 2020 é de \qty{500}{Gt} de dióxido de carbono com uma probabilidade de \qty{50}{\percent}, e
\qty{1350}{Gt} para \qty{2}{\celsius}. As escolhas feitas sobre os outros fatores deslocam esses valores em
\qty{220}{Gt} ou mais; o problema de fim de semana reconstrói o primeiro.
% ledger: ar6ch5:budget-parts, ar6spm:budget

\begin{dated}{2025}
% ledger: data:co2-mlo, data:gistemp (values read from the datasets as fetched 2026-10-22; offset computed)
A média anual do dióxido de carbono em Mauna Loa atingiu \qty{427.35}{ppm} em 2025, um forçamento de
$3.93\log_2(427.35/278) = \qty{2.44}{W.m^{-2}}$ pela lei logarítmica. No registro de temperatura
da NASA, 2024 é o ano mais quente, \qty{1.50}{\celsius} acima da média de 1880--1900, os primeiros
anos desse registro; 2025 ficou em \qty{1.40}{\celsius}.
\end{dated}

\begin{omfigure}
\centering
\begin{tikzpicture}
\begin{axis}[omrecord, width=0.86\linewidth, height=5.8cm, xmin=0, xmax=4000, ymin=0, ymax=2.6,
  xlabel={emissões cumulativas de \ce{CO2} desde 1850 (\unit{Gt})},
  ylabel={aquecimento desde 1850--1900 (\unit{\celsius})}, xtick={0,500,...,4000},
  /pgf/number format/1000 sep={\,}, xtick pos=bottom, ytick pos=left, clip=false,
  tick label style={font=\scriptsize}, label style={font=\small}]
% computed in figdata/bachelor-2/anthropogenic.py part tcre; ledger: ar6spm:tcre, ar6spm:budget
\addplot[name path=hi, draw=none, forget plot] table[x=E_GtCO2, y=dT_hi] {figdata/out/bachelor-2-anthropogenic-tcre.dat};
\addplot[name path=lo, draw=none, forget plot] table[x=E_GtCO2, y=dT_lo] {figdata/out/bachelor-2-anthropogenic-tcre.dat};
\addplot[omThm!18, forget plot] fill between[of=hi and lo];
\addplot[omThm, line width=1.3pt] table[x=E_GtCO2, y=dT_mid] {figdata/out/bachelor-2-anthropogenic-tcre.dat};
\draw[densely dashed, black!60] (axis cs:0,1.5) -- (axis cs:4000,1.5);
\draw[densely dashed, black!60] (axis cs:0,2.0) -- (axis cs:4000,2.0);
\node[anchor=north west, font=\scriptsize, text=black!70] at (axis cs:0,1.5) {\qty{1.5}{\celsius}};
\node[anchor=north west, font=\scriptsize, text=black!70] at (axis cs:0,2.0) {\qty{2}{\celsius}};
\fill[black] (axis cs:2390,1.07) circle (2.2pt);
\draw[black, line width=0.8pt] (axis cs:2150,1.07) -- (axis cs:2630,1.07);
\draw[black, line width=0.8pt] (axis cs:2390,0.8) -- (axis cs:2390,1.3);
\node[anchor=north west, font=\scriptsize, align=left] at (axis cs:2440,1.02) {1850--2019:\\\qty{2390}{Gt}, \qty{1.07}{\celsius}};
\draw[<->, omDef, thick] (axis cs:2390,1.62) -- (axis cs:2890,1.62);
\node[anchor=south, font=\scriptsize, text=omDef, align=center] at (axis cs:2640,1.64) {AR6: \qty{500}{Gt}};
\node[anchor=north west, font=\scriptsize, text=omThm, align=left] at (axis cs:150,2.5) {intervalo provável da TCRE,\\0.27 a \qty{0.63}{\celsius} por \qty{1000}{Gt}};
\end{axis}
\end{tikzpicture}
\par\smallskip{\small Aquecimento em função das emissões cumulativas de dióxido de carbono para a TCRE avaliada,
\qty{0.45}{\celsius} por \qty{1000}{Gt} (linha) e seu intervalo provável (sombreado), com o aquecimento humano
e as emissões até hoje (ponto, intervalos prováveis como barras). A seta marca o orçamento restante para
\qty{1.5}{\celsius} a \qty{50}{\percent} a partir de 2020: menor do que uma leitura da reta sozinha,
porque os outros fatores também aquecem.}
\end{omfigure}


\section{O nível do mar}\label{sec:b2:anthropogenic-climate-change:sea-level}

\begin{definition}[Balanço do nível do mar; elevação termostérica do nível do mar]\label{def:b2:anthropogenic-climate-change:sea-level}
O \emph{balanço do nível do mar}\index{balanço do nível do mar} é a comparação da variação observada do
nível médio global do mar com a soma de suas contribuições estimadas: a dilatação do oceano que
aquece, o derretimento das geleiras e dos mantos de gelo da Groenlândia e da Antártida, e a variação da
água armazenada nos continentes. A \emph{elevação termostérica do nível do mar}\index{elevação termostérica do nível do mar} é a
parte devida à dilatação térmica da água do mar, com massa constante.
\end{definition}

\begin{proposition}[Dilatação térmica de uma camada]\label{prop:b2:anthropogenic-climate-change:thermosteric}
Uma camada de oceano de espessura $H$, com coeficiente de dilatação térmica $\alpha$, que aquece uniformemente
de $\Delta T$ com massa e área constantes sobe $\Delta h = \alpha\,\Delta T\,H$.
\end{proposition}
\begin{proof}
A massa $\rho H A$ de uma coluna de área $A$ é constante, logo $\Delta H/H = -\Delta\rho/\rho =
\alpha\,\Delta T$ em primeira ordem, com $\alpha = -(1/\rho)\,\partial\rho/\partial T$.
\end{proof}
% ledger: teos10:alpha-beta

No oceano real $\alpha$ cresce com a temperatura e a pressão, e o aquecimento concentra-se perto
da superfície. A avaliação resume o resultado como uma eficiência de dilatação de \qty{0.113}{m} por
\qty{e24}{J} de calor acrescentado ao oceano, ou \qty{0.617}{m} por grau de aquecimento do oceano inteiro
(\emph{média confiança}). Sobre sua profundidade média de \qty{3682}{m}, isso é um $\alpha$ efetivo de
$0.617/3682 = \qty{1.68e-4}{\celsius^{-1}}$ (\cref{exo:b2:anthropogenic-climate-change:8}).
% ledger: ar6ch9:expansion, charette:ocean-volume

O balanço, desenhado na figura abaixo, fecha dentro de suas incertezas. Em 1971--2018 as contribuições somam
\qty{2.00}{mm.yr^{-1}}, contra uma subida observada de \qty{2.33}{mm.yr^{-1}}; em 2006--2018 elas
somam 3.61 contra 3.69. A dilatação térmica foi o maior termo isolado, metade da soma em 1971--2018, mas o gelo está
alcançando-a: em 2006--2018 as geleiras e os mantos de gelo juntos deram \qty{1.62}{mm.yr^{-1}}, mais que
a dilatação. O oceano também continua a aquecer em profundidade muito depois que a superfície se estabiliza, e os mantos
de gelo respondem ao longo de séculos; o nível do mar está comprometido a subir durante séculos a milênios e
permanecerá alto por milhares de anos (\emph{alta confiança}). Nos próximos dois mil anos,
projeta-se que suba cerca de 2 a \qty{3}{m} se o aquecimento ficar em \qty{1.5}{\celsius}, 2 a
\qty{6}{m} a \qty{2}{\celsius} e 19 a \qty{22}{m} a \qty{5}{\celsius} (\emph{baixa confiança}).
% ledger: ar6ch9:table95, ar6spm:sea-commit

\begin{omfigure}
\centering
\begin{tikzpicture}
\begin{axis}[ybar stacked, width=0.72\linewidth, height=6.2cm, ymin=0, ymax=4.3, bar width=26pt,
  xtick=data, symbolic x coords={1971--2018, 1993--2018, 2006--2018}, enlarge x limits=0.25,
  ylabel={taxa de subida (\unit{mm.yr^{-1}})}, tick label style={font=\scriptsize},
  label style={font=\small}, ytick pos=left, xtick pos=bottom,
  legend style={at={(1.03,0.5)}, anchor=west, font=\scriptsize, draw=black!30},
  legend cell align=left, reverse legend]
% ledger: ar6ch9:table95
\addplot[fill=omDef!70, draw=omDef!90!black] coordinates {(1971--2018,1.01) (1993--2018,1.31) (2006--2018,1.39)};
\addlegendentry{dilatação térmica}
\addplot[fill=omDef!30, draw=omDef!70!black] coordinates {(1971--2018,0.44) (1993--2018,0.55) (2006--2018,0.62)};
\addlegendentry{geleiras}
\addplot[fill=omThm!55, draw=omThm!80!black] coordinates {(1971--2018,0.25) (1993--2018,0.43) (2006--2018,0.63)};
\addlegendentry{Groenlândia}
\addplot[fill=omThm!25, draw=omThm!70!black] coordinates {(1971--2018,0.14) (1993--2018,0.25) (2006--2018,0.37)};
\addlegendentry{Antártida}
\addplot[fill=black!20, draw=black!50] coordinates {(1971--2018,0.15) (1993--2018,0.31) (2006--2018,0.60)};
\addlegendentry{água continental}
\node[fill=black, rotate=45, minimum size=5pt, inner sep=0pt] at (axis cs:1971--2018,2.33) {};
\node[fill=black, rotate=45, minimum size=5pt, inner sep=0pt] at (axis cs:1993--2018,3.25) {};
\node[fill=black, rotate=45, minimum size=5pt, inner sep=0pt] at (axis cs:2006--2018,3.69) {};
\end{axis}
\end{tikzpicture}
\par\smallskip{\small O \omterm{def:b2:anthropogenic-climate-change:sea-level}{balanço do nível do mar}: contribuições observadas para a subida do nível médio global
do mar, como taxas em três períodos que terminam em 2018, empilhadas, e a taxa observada (losangos). A
diferença entre a pilha e o losango é a parte que o balanço não fecha. Valores: Sexto Relatório de
Avaliação do IPCC, Grupo de Trabalho~I, Tabela 9.5.}
\end{omfigure}

\section*{Exercícios}\label{sec:b2:anthropogenic-climate-change:exercises}

Dados para os exercícios, salvo indicação em contrário: forçamento do dióxido de carbono duplicado $F_{2\times} =
\qty{3.93}{W.m^{-2}}$ a partir de $C_0 = \qty{278}{ppm}$; constante de Stefan--Boltzmann $\sigma =
\qty{5.670e-8}{W.m^{-2}.K^{-4}}$; raio da Terra \qty{6371}{km}; um ano
\qty{3.156e7}{s}; massas molares do carbono e do dióxido de carbono 12.011 e \qty{44.009}{g.mol^{-1}};
as retroalimentações e o \omterm{def:b2:anthropogenic-climate-change:sea-level}{balanço do nível do mar} do capítulo.
% ledger: ipcc:erf-2xco2, const:sigma, earth:Rmean, ciaaw:abridged-2024, ar6ch7:table710, ar6ch9:table95

\begin{exercise}[$\star$]\label{exo:b2:anthropogenic-climate-change:1}
O forçamento do dióxido de carbono duplicado é de \qty{3.75}{W.m^{-2}} depois do ajuste da estratosfera;
a troposfera então se ajusta em $-0.60$ (temperatura), +0.22 (vapor d'água), +0.45 (nuvens) e
+0.11 (superfície terrestre) \unit{W.m^{-2}}. (a) Encontre o ERF. (b) Em que porcentagem ele excede os
\qty{3.71}{W.m^{-2}} da avaliação anterior? (c) Qual ajuste é negativo, e por que uma troposfera mais quente
reduz o forçamento?
\end{exercise}
% ledger: ar6ch7:2xco2

\begin{exercise}[$\star$]\label{exo:b2:anthropogenic-climate-change:2}
O sistema climático ganhou \qty{434.9}{ZJ} de 1971 a 2018 e \qty{152.4}{ZJ} de 2006 a 2018,
contados entre os meios do primeiro e do último ano. (a) Encontre a área da superfície da Terra.
(b) Converta cada ganho num desequilíbrio médio em \unit{W.m^{-2}}. (c) O oceano recebeu
\qty{91}{\percent}: quantos zettajoules em 1971--2018?
\end{exercise}
% ledger: ar6ch7:inventory, earth:Rmean

\begin{exercise}[$\star$]\label{exo:b2:anthropogenic-climate-change:3}
(a) Some os valores centrais das retroalimentações da figura, e compare com a líquida. (b) Encontre
a ECS com o valor líquido, e o aquecimento com a \omterm{def:b2:anthropogenic-climate-change:feedbacks}{resposta de Planck} sozinha. (c) Por que fator
as outras retroalimentações amplificam o aquecimento de Planck?
\end{exercise}

\begin{exercise}[$\star$]\label{exo:b2:anthropogenic-climate-change:4}
(a) Converta uma TCRE de \qty{0.45}{\celsius} por \qty{1000}{Gt} de dióxido de carbono em graus por
\qty{1000}{PgC} de carbono, e faça o mesmo para seu intervalo provável, 0.27 a 0.63. (b) Converta os
\qty{2390}{Gt} de dióxido de carbono emitidos de 1850 a 2019 em petagramas de carbono. (c) Que
aquecimento a TCRE dá para essas emissões?
\end{exercise}
% ledger: ar6spm:tcre, ar6spm:budget

\begin{exercise}[$\star\star$]\label{exo:b2:anthropogenic-climate-change:5}
(a) A superfície emite \qty{398}{W.m^{-2}} e o efeito estufa retém \qty{159}{W.m^{-2}} desse
fluxo. Encontre $g$ e a temperatura da superfície de um corpo negro que emite \qty{398}{W.m^{-2}}. (b) Encontre a
\omterm{def:b2:anthropogenic-climate-change:feedbacks}{resposta de Planck} a \qty{288}{K}, e a \qty{298}{K} com o mesmo $g$. (c) A \omterm{def:b2:anthropogenic-climate-change:feedbacks}{resposta de Planck}
sozinha tornaria uma Terra mais quente mais ou menos sensível?
\end{exercise}
% ledger: ipcc:budget-surface, ar6ch7:planck

\begin{exercise}[$\star\star$]\label{exo:b2:anthropogenic-climate-change:6}
Os satélites dão uma \omterm{def:b2:anthropogenic-climate-change:feedbacks}{retroalimentação do vapor d'água} de $1.85 \pm 0.32$ e os estudos uma \omterm{def:b2:anthropogenic-climate-change:feedbacks}{retroalimentação do gradiente vertical} de
$-0.50 \pm 0.20$ \unit{W.m^{-2}.\celsius^{-1}} (um desvio-padrão). (a) Encontre sua soma.
(b) Se as duas fossem independentes, qual seria o desvio-padrão da soma? (c) Os modelos
dão à soma um desvio-padrão de apenas \num{0.10}. O que isso diz sobre os dois erros?
\end{exercise}
% ledger: ar6ch7:wv-lr

\begin{exercise}[$\star\star$]\label{exo:b2:anthropogenic-climate-change:7}
De 1850--1900 a 2006--2019, $\Delta N = \qty{0.59}{W.m^{-2}}$, $\Delta T = \qty{1.03}{\celsius}$
e $\Delta F = \qty{2.20}{W.m^{-2}}$. (a) Encontre $\lambda$ e a sensibilidade efetiva. (b) Espera-se que o
parâmetro de retroalimentação se torne menos negativo em $\lambda' = \qty{0.5}{W.m^{-2}.\celsius^{-1}}$
à medida que o padrão de aquecimento evolui. Encontre a sensibilidade $-F_{2\times}/(\lambda + \lambda')$.
(c) Refaça (b) com $\lambda' = 0.1$.
\end{exercise}
% ledger: ar6ch7:energy-budget, ar6ch7:pattern, ar6ch7:pattern-ecs

\begin{exercise}[$\star\star$]\label{exo:b2:anthropogenic-climate-change:8}
(a) A partir dos \qty{0.617}{m} de subida por grau de aquecimento do oceano inteiro, profundidade média
\qty{3682}{m}, encontre o coeficiente de dilatação efetivo. (b) O oceano ganhou \qty{396}{ZJ} em
1971--2018. Com uma eficiência de dilatação de \qty{0.113}{m} por \qty{e24}{J}, encontre a subida
termostérica, e compare com os \qty{47.5}{mm} do balanço. (c) Uma camada de \qty{700}{m} aquece
\qty{0.2}{\celsius} com $\alpha = \qty{2e-4}{\celsius^{-1}}$: quanto ela eleva o nível do mar?
\end{exercise}
% ledger: ar6ch9:expansion, charette:ocean-volume, ar6ch7:inventory, ar6ch9:table95

\begin{exercise}[$\star\star$]\label{exo:b2:anthropogenic-climate-change:9}
Uma impressão digital de modelo, em quatro décadas, é $x = (0.10, 0.20, 0.30, 0.40)$ \unit{\celsius}; as
observações são $y = (0.12, 0.18, 0.33, 0.41)$ \unit{\celsius}. (a) Encontre o fator de escala de mínimos quadrados
$\beta$ de $y = \beta x + \varepsilon$. (b) Seu erro-padrão, estimado a partir de simulações sem
forçamento, é \num{0.08}. A resposta é detectada? É coerente com o modelo? (c) O que
significaria um $\beta$ de 0.5 com o mesmo erro?
\end{exercise}

\begin{exercise}[$\star\star\star$]\label{exo:b2:anthropogenic-climate-change:10}
Duas impressões digitais, gases de efeito estufa $g = (0.2, 0.4, 0.6, 0.8)$ e aerossóis $a = (-0.1, -0.2,
-0.2, -0.3)$, ambas em \unit{\celsius} ao longo de quatro décadas, são ajustadas às observações $y = (0.1,
0.2, 0.4, 0.5)$. (a) Escreva as equações normais do ajuste de mínimos quadrados $y = \beta_g g + \beta_a a
+ \varepsilon$ e resolva-as. (b) Encontre o aquecimento atribuído a cada fator na última década.
(c) Por que as incertezas sobre $\beta_g$ e $\beta_a$ crescem quando as duas impressões digitais se parecem?
\end{exercise}

\begin{exercise}[$\star\star\star$]\label{exo:b2:anthropogenic-climate-change:11}
Use o método do \omterm{def:b2:anthropogenic-climate-change:budget}{orçamento de carbono restante} com a TCRE de \qty{0.45}{\celsius} por
\qty{1000}{Gt}, os \qty{210}{Gt} emitidos de 2015 a 2019 e \qty{26}{Gt} por grau para as
retroalimentações do sistema terrestre. (a) Encontre o orçamento para \qty{2}{\celsius}, \qty{0.93}{\celsius} acima de
2010--2019, se os outros fatores acrescentarem \qty{0.1}{\celsius}. (b) Que aquecimento adicional pelos outros
fatores o torna igual aos \qty{1350}{Gt} avaliados? Isso está dentro do intervalo avaliado? (c) Mostre
que cada \qty{0.1}{\celsius} de aquecimento não devido ao CO$_2$ altera qualquer orçamento da mesma quantidade.
\end{exercise}

\begin{exercise}[$\star\star\star$]\label{exo:b2:anthropogenic-climate-change:12}
A retroalimentação líquida tem um intervalo muito provável de $-1.81$ a \qty{-0.51}{W.m^{-2}.\celsius^{-1}}, em torno de
$-1.16$. (a) Encontre a ECS nas duas extremidades e no centro. (b) Mostre que um intervalo de $\lambda$
simétrico em torno de seu centro dá um intervalo de ECS que não o é, e diga qual lado é mais longo. (c) Por que
as avaliações acham mais difícil excluir uma ECS alta que uma baixa?
\end{exercise}
% ledger: ar6ch7:table710

\begin{problem}[{Problema de fim de semana --- ler o balanço}]\label{pb:b2:anthropogenic-climate-change:1}
Este problema segue a energia e o carbono da avaliação desde o forçamento até o orçamento
restante para \qty{1.5}{\celsius}. Use os dados dos exercícios, e: ERF 1750--2019 (\unit{W.m^{-2}})
dióxido de carbono 2.16, metano 0.54, óxido nitroso 0.21, gases halogenados 0.41, ozônio 0.47,
vapor d'água estratosférico 0.05, aerossóis--radiação $-0.22$, aerossóis--nuvens $-0.84$, uso da terra
$-0.20$, carbono negro sobre a neve 0.08, rastros de condensação 0.06; dióxido de carbono \qty{410}{ppm} em 2019 e
\qty{427.35}{ppm} em 2025; taxa de aquecimento \qty{0.79}{W.m^{-2}} em 2006--2018; as variações do balanço de energia
$\Delta N = 0.59$, $\Delta F = \qty{2.20}{W.m^{-2}}$, $\Delta T = \qty{1.03}{\celsius}$; TCRE
\qty{0.45}{\celsius} por \qty{1000}{Gt} de dióxido de carbono (provável 0.27 a 0.63); \qty{2390}{Gt}
emitidos em 1850--2019; aquecimento humano \qty{1.07}{\celsius} para 2010--2019; \qty{210}{Gt} emitidos de
2015 a 2019; aquecimento adicional não devido ao CO$_2$ \qty{0.1}{\celsius}, o limite inferior dos 0.1
a \qty{0.2}{\celsius} avaliados; retroalimentações do sistema terrestre \qty{26}{Gt} por grau. Arredonde a três algarismos
significativos.
% ledger: ipcc:erf-agents, ipcc:erf-2xco2, data:co2-mlo, ar6ch7:inventory, ar6ch7:energy-budget, ar6spm:tcre, ar6spm:budget, ar6ch5:budget-parts

\noindent\textbf{Parte I --- O forçamento.}
\begin{enumerate}
  \item Some o forçamento dos quatro gases de efeito estufa bem misturados.
  \item Some o forçamento dos aerossóis.
  \item Some todos os termos humanos. Compare com os \qty{2.72}{W.m^{-2}} avaliados.
  \item Verifique o forçamento do dióxido de carbono em 2019 com a lei logarítmica.
  \item Encontre o forçamento do dióxido de carbono em 2025 pela mesma lei.
  \item Qual termo tem a maior incerteza, e por que ele domina a do total?
\end{enumerate}

\noindent\textbf{Parte II --- O balanço de energia.}
\begin{enumerate}[resume]
  \item Quanta energia o sistema climático ganhou em 2006--2018, em zettajoules?
  \item Quanto dela foi para o oceano?
  \item Encontre o parâmetro de retroalimentação a partir das variações do balanço de energia.
  \item Encontre a \omterm{def:b2:anthropogenic-climate-change:ecs}{sensibilidade climática de equilíbrio} efetiva.
  \item Encontre a \omterm{def:b2:anthropogenic-climate-change:tcr}{resposta climática transitória} a partir dos mesmos dados.
  \item Encontre a ECS a partir da retroalimentação líquida da figura, e compare os dois valores.
  \item Dê a razão física pela qual a estimativa histórica é a mais baixa.
\end{enumerate}

\noindent\textbf{Parte III --- Aquecimento por tonelada.}
\begin{enumerate}[resume]
  \item Converta a TCRE em graus por \qty{1000}{PgC} de carbono.
  \item Converta seu intervalo provável.
  \item Converta as emissões de 1850--2019 em petagramas de carbono.
  \item Que aquecimento a TCRE dá para as emissões de 1850--2019?
  \item Compare com o aquecimento humano de \qty{1.07}{\celsius}. O que a comparação diz sobre o
        efeito líquido dos outros fatores humanos?
\end{enumerate}

\noindent\textbf{Parte IV --- O orçamento restante para \qty{1.5}{\celsius}.}
\begin{enumerate}[resume]
  \item Que aquecimento resta acima de 2010--2019?
  \item Retire o aquecimento não devido ao CO$_2$ e divida pela TCRE: que orçamento resta a partir do meio
        de 2010--2019?
  \item Retire as emissões de 2015--2019.
  \item Retire as retroalimentações do sistema terrestre sobre o aquecimento restante.
  \item Em quanto muda o orçamento se os outros fatores acrescentarem \qty{0.2}{\celsius}?
  \item Em resumo: qual é o \omterm{def:b2:anthropogenic-climate-change:budget}{orçamento de carbono restante} para \qty{1.5}{\celsius}, a partir do início de
        2020, calculado a partir da TCRE? Compare com os \qty{500}{Gt} avaliados.
\end{enumerate}
\end{problem}
